\section{Additional social context and counterexamples}
\label{app:social-detail}

\paragraph{G11: a transient echo, a check and a rejection.}
At \decision{11}{255} Sol (seat~4) offers Virginia for Grok's (seat~2) mortgaged Pennsylvania plus \$150, publicly asserting ``Virginia completes your pinks''. Grok holds States but not St.~Charles, which Sol owns, so the claim is false. Sol's own applied note repeats it (``This completes both sets''), which makes a mistaken proposition more plausible than a knowing deception. Grok's returned reasoning at d256 first echoes the premise (``accepting would complete my pink monopoly'') and, a few hundred characters later, corrects it (``would not complete''). Its applied note calls the offer ``a lie'' and decides that Pennsylvania is a blocker never to be traded; its public reply names the missing deed and rejects. Sol concedes at d257 (``You're right---I missed St.~Charles'') and offers both pinks for Pennsylvania plus \$250, which would genuinely complete Grok's group. Grok rejects that too at d258, stating that the deal would complete Sol's greens. An early textual echo is not the final response; a recipient's accusation of lying does not establish the proposer's prior knowledge; and own completion is not sufficient for acceptance in this case.

\paragraph{G4: a public concession with prior correct knowledge.}
After Sol (seat~3) completes pinks, Gemini (seat~4) publicly invites a deal so Sol is not the only player with a monopoly at d138. Grok (seat~1) answers at d139, then buys mortgaged Ventnor from Sol for \$200 at d146 to make a yellow set tradable. Its notes weigh how many opposing pink houses that payment could fund. In the orange--yellow negotiation, Gemini's d148 applied note records the true \$150 yellow house cost and a contingency: if Grok insists on orange, Gemini can demand compensation. Grok's d149 message falsely says yellow costs \$200 per house. Gemini's d150 public ``Fair point on house costs'' precedes a request to take yellow with \$143, but its returned explanation still says \$150, not \$200. Grok counters at \$50 and Gemini accepts at d152. Events transfer Marvin, mortgaged Ventnor, Water Works and \$50 to Gemini for St.~James and New York, charging Gemini \$13 mortgage interest. Both obtain complete groups. Gemini redeems Ventnor at d175 and builds at d176. The direction change matches a contingency recorded before the false figure, leaving any price or framing effect unidentified.

\paragraph{G18: denial that ends with creditor inheritance.}
Grok (seat~3) buys mortgaged Illinois from Gemini for \$100 at d335--336, saying this keeps Sol (seat~1) off the reds. Eighteen turns later, at d360, Grok owes Sol \$950 with \$64 cash and declares bankruptcy. Illinois and both dark blues pass to Sol; Illinois completes the reds the purchase was meant to block. The observed denial lasts only as long as its holder remains solvent. This parallels G13's later Short Line estate transfer in Transcript~\ref{tr:omission}. Neither observation proves denial or solvency purchases valueless, and the debt fixes the creditor rather than allowing the bankrupt player to select a beneficiary.

A separate G18 record completes the eventual winner's first group through a received-call conflict. At d83 Grok's received call accepts \$130 for Baltic, completing Opus's browns. Its assistant content selects rejection and its applied note says to keep Baltic as a block. This has a material immediate consequence, but the contribution to Opus's eventual win is untested. It does not support an intended-concession interpretation.

\paragraph{Other refusals to fund rivals.}
Gemini's G9 d82 rejection says that paying would fund a dark-blue house, and G16 d322 says keeping the rival cash-starved is preferable. These are stated reasons; their rival-cash mechanisms were not recomputed in the writing-stage check. They complement the original-checked G20 appeal in Transcript~\ref{tr:boycott}, without supplying a prevalence denominator. The review did not code every trade response for leader-funding reasons. No multi-party pact, cartel or deliberate kingmaking is identified in these adjudicated examples; this is not a corpus-wide absence estimate. Designed collusion studies specify additional opportunities and incentives \citep{shi2026emergentcollusion}.
